%% 02_literature.tex
%% Section 2: Literature Review

\section{Literature Review}
\label{sec:literature}

This paper draws on and contributes to three intersecting literatures: the economics of stranded assets, climate risk and financial stability, and Indonesia-specific research on coal transition risk. We review each in turn before identifying the specific gap that our integrated framework addresses.

\subsection{The Stranded Assets Literature}
\label{subsec:stranded_assets}

The foundations of the stranded assets concept can be traced to the carbon budget literature, which establishes an upper bound on cumulative greenhouse gas emissions compatible with specific temperature targets. \citet{McGlade2015} translate this carbon budget constraint into geographical and fuel-specific terms, estimating that 82\% of global coal reserves, 49\% of gas reserves, and 33\% of oil reserves must remain unextracted to maintain a 50\% probability of limiting warming to 2\textdegree C. \citet{Welsby2021} update this analysis for the more ambitious 1.5\textdegree C target, finding that 89\% of coal reserves globally are ``unextractable,'' with particularly severe implications for coal-exporting nations in Southeast Asia and Africa.

The translation from physical unextractability to economic value destruction is the domain of stranded asset analysis proper. \citet{CarbonTracker2012} used the term through their ``carbon bubble'' thesis, arguing that fossil fuel companies are systematically overvalued because markets fail to price the risk that reserves will become economically unviable. \citet{Caldecott2021} provide a comprehensive taxonomy, distinguishing between demand-side stranding (declining coal consumption due to renewable energy competition and climate policy), supply-side stranding (regulatory restrictions on extraction), and cost-side stranding (rising extraction costs as marginal deposits are developed). \citet{VanDerPloeg2020} formalise these dynamics in a resource economics framework, showing that the announcement of future climate policy can trigger immediate asset revaluation as rational agents update expectations.

Empirical quantification of stranded values has proceeded along two tracks. The first employs supply cost curve analysis, ranking fossil fuel deposits by extraction cost and determining which become uneconomic as demand declines or carbon prices rise. \citet{Mendelevitch2019} applies this approach to global steam coal, demonstrating that Indonesian coal is particularly vulnerable owing to its high transport costs to distant markets and competition from lower-cost Australian producers. The second track uses discounted cash flow (DCF) models to value individual firms or projects under alternative climate scenarios. \citet{Edwards2022} apply this methodology to coal-fired power plants, estimating committed emissions and the capital at risk from early retirement. \citet{Pfeiffer2018} extend the analysis to planned infrastructure, finding that committed emissions from existing and planned power plants alone could exhaust the 1.5\textdegree C carbon budget.

More recent work has sharpened both the physical and economic case for coal stranding. \citet{Tong2019} estimate that existing and planned coal-fired power plants alone would exhaust the remaining 1.5\textdegree C carbon budget if operated over their full design lifetimes, implying a global coal phase-out that is not merely a long-run projection but a near-term imperative. \citet{Cui2019} construct a mine-level global coal supply curve to identify which deposits are unburnable under a 1.5\textdegree C pathway, finding that Southeast Asian sub-bituminous coal, including Indonesian deposits, is among the most cost-disadvantaged given its low calorific value relative to its transport cost to Asian markets. Together, these studies reinforce the urgency of Indonesia-specific stranded asset quantification.

Our work is closely related to \citet{sonny2016}, who provide the first Indonesia-specific quantification using a cash cost approach applied to 16 listed companies, estimating aggregate stranded values of \$58--99 billion depending on discount rate and mine life assumptions (central estimate: \$86 billion at 8\% discount, 20-year life, 40\% price decline). While their firm-level estimates and fiscal discussion are valuable, their analysis does not extend to the banking sector or incorporate scenario-contingent DCF modelling with Monte Carlo uncertainty quantification.

All in all, a general limitation of this literature is its predominantly physical or sectoral focus. Most studies quantify the volume of unburnable reserves or the value at risk within the fossil fuel sector, but do not trace how these losses propagate to the financial system or the broader economy. This paper addresses that limitation by linking firm-level stranding estimates to bank exposures and macroeconomic outcomes.

\subsection{Climate Risk and Financial Stability}
\label{subsec:climate_risk}

The recognition that climate change poses risks to financial stability has accelerated rapidly since Mark Carney's influential ``tragedy of the horizon'' speech in 2015 and the establishment of the Network for Greening the Financial System (NGFS) in 2017. \citet{Battiston2017} provide one of the first quantitative assessments, stress-testing the European financial system against a sudden fossil fuel price decline and finding that direct exposures to carbon-intensive sectors represent a material but manageable risk for most European banks. However, they note that indirect exposures through interbank lending and derivatives could amplify losses substantially.

Subsequent climate stress tests have been conducted by the Dutch central bank \citep{Vermeulen2018}, the Bank of England \citep{BoE2019}, and the European Central Bank \citep{ECB2022}. These exercises reveal several common findings: transition risk is concentrated in specific sectors (energy, transport, heavy industry); bank losses depend critically on whether the transition is orderly or disorderly; and second-round effects, including fire sales, counterparty defaults, and funding market disruptions, can substantially amplify first-round losses. The IMF has extended this framework to emerging markets, including a Financial Sector Assessment Program for Indonesia \citep{IMF2024Indonesia} that identifies coal exposure as a material risk but stops short of a full climate stress test.

On the asset pricing side, \citet{BoltonKacperczyk2021} document a carbon risk premium in equity markets, showing that firms with higher carbon emissions earn higher stock returns, consistent with investors demanding compensation for transition risk. \citet{Dietz2016} estimate the climate value-at-risk of global financial assets at \$2.5 trillion under a business-as-usual scenario, with a fat-tailed distribution that implies non-trivial probability of losses exceeding \$24 trillion. \citet{Daumas2024} provides a comprehensive survey of this literature, noting that while theoretical understanding of climate-financial linkages has advanced considerably, empirical evidence grounded in actual firm-level data remains scarce, particularly for emerging markets.

The methodological frontier in this literature has moved toward network-based stress testing approaches. \citet{Roncoroni2021} demonstrate that bilateral financial exposures between institutions can generate amplification dynamics in which second-round effects substantially exceed first-round losses, particularly in concentrated networks with high overlap between lenders. \citet{Battiston2021} develop a climate policy-relevant extension of systemic risk analysis, showing that transition risk shocks propagate through both direct exposures and indirect network channels in ways that aggregate stress testing misses. For emerging markets specifically, \citet{IMF2024Indonesia} note that the combination of bank-dominated financial systems, concentrated industrial structures, and limited capital market alternatives creates a particularly fertile ground for credit channel amplification, yet standard stress testing tools applied to advanced economies cannot be mechanically transplanted to this context.

Our contribution to this literature is twofold. First, we construct a bipartite bank-coal exposure network from actual loan data rather than relying on sectoral aggregates, allowing us to identify specific bank-firm channels of risk transmission. Second, we conduct what is, to our knowledge, the first bottom-up climate stress test of the Indonesian banking system that traces losses from individual coal company stranding through identified credit exposures to bank-level capital adequacy impacts.

\subsection{Indonesia: Coal, Banking, and the Transition Gap}
\label{subsec:indonesia_gap}

Indonesia-specific research on climate transition risk has grown in recent years but remains fragmented. \citet{IESR2023} provides the most comprehensive assessment of coal as stranded assets in Indonesia, using a combination of reserve analysis, scenario modelling, and qualitative banking assessment. Their report identifies 13 of 34 coal companies as highly vulnerable to stranding under a 1.5\textdegree C scenario and notes that four major banks hold significant coal exposure, but does not quantify credit losses or macroeconomic impacts. The IMF's Financial Sector Assessment Program for Indonesia \citep{IMF2024Indonesia, IMF2025IndonesiaStress} incorporates climate risk as a supplementary consideration but focuses primarily on physical risks (flooding, sea-level rise) rather than transition risk from coal stranding.

\citet{Rishanty2023}, in a Bank Indonesia working paper, estimate the total climate transition value-at-risk at \$206 billion (NPV, 2023) under a Well Below 2\textdegree C scenario, distributed across investors (\$35 billion), government (\$51 billion), and the coal supply chain (\$120 billion). While they emphasise the role of banks in financing local transitions, they do not construct a bank-level stress test or quantify credit losses through identified exposure channels.

\citet{Burke2018} examine Indonesia's energy policy landscape, documenting the deep structural dependence on coal for both electricity generation and export revenues. \citet{Makhijani2015} quantify fossil fuel subsidies in Indonesia, which further entrench coal dependence by artificially lowering the cost of coal-fired electricity relative to renewables. \citet{JakobSteckel2016} situate Indonesia within the broader political economy of energy transitions in developing countries, noting that concentrated coal interests and dispersed transition benefits create formidable barriers to policy change.

To summarise, while prior studies have quantified the magnitude of Indonesian coal stranded assets at the firm level \citep{sonny2016} or the mine level \citep{Rishanty2023}, and have qualitatively noted the banking sector implications \citep{IESR2023}, no existing study provides an integrated, quantitative assessment of how coal stranding transmits through the banking sector via identified credit exposures. Table~\ref{tab:comparison} in Section~\ref{sec:discussion} provides a detailed comparison of our approach with prior work, highlighting the novelty of our firm-level, bottom-up, two-model framework.

\subsection{Research Gap and Our Contribution}
\label{subsec:research_gap}

The three literatures reviewed above converge on a shared lacuna: while much is known about the \textit{physical} extent of coal stranding risk globally, about the \textit{methodological} tools for climate stress testing in advanced financial systems, and about the \textit{political economy} of Indonesia's coal dependence, no study integrates these strands into a coherent, empirically grounded framework for a single coal-dependent emerging economy.

Specifically, five gaps motivate the present paper. First, existing stranded asset quantifications for Indonesia \citep{sonny2016, Rishanty2023} do not extend to bank exposures: they estimate coal value losses but stop at the boundaries of the mining sector. Second, existing climate stress tests for Indonesia are either qualitative \citep{IESR2023} or treat the coal sector as one component of a broader physical-risk assessment \citep{IMF2024Indonesia}, rather than conducting a dedicated bottom-up stress test of transition risk through identified credit channels. Third, the bipartite network structure of bank-coal lending in Indonesia, which determines the distribution of losses across financial institutions and creates potential for correlated defaults, has not been systematically documented. Fourth, the literature on the sovereign-bank nexus \citep{Arellano2008} and on fiscal vulnerability in resource-dependent economies \citep{VanDerPloegVenables2011} has not been applied to the coal transition context, where the interaction between fiscal revenue losses and banking sector stress is potentially self-reinforcing. Fifth, no existing study provides the firm-level Monte Carlo uncertainty quantification necessary to move beyond deterministic point estimates and construct probability-weighted, confidence-interval-bounded stranded value assessments.

Our paper addresses each of these gaps. We contribute: (i) the first granular, scenario-contingent, Monte Carlo-bounded estimate of stranded coal asset values for listed Indonesian companies; (ii) the first bipartite bank-coal exposure network constructed from loan-level disclosures; (iii) the first bank-level climate stress test for Indonesia calibrated to this network; and (iv) a clear documentation of the coal-banking-sovereign nexus in its Indonesian manifestation. Together, these contributions provide a replicable methodology and a concrete empirical baseline for regulators, investors, and researchers working on transition risk in coal-dependent emerging economies.
